\section{Discussion: admissibility, support, and residual interpretation}
\subsection{Storage coordinates need not be terminal-current coordinates}
The real-data study makes a distinction necessary: a conservative magnetic
subsystem does not automatically imply a potential whose gradient against
measured terminal current is the reconstructed flux. A schematic loss branch
has $\statorcurrent=\magnetizingcurrent+\vect i_{\rm Fe}$.
The storage relation can be conservative in $\magnetizingcurrent$ while its
stationary reduction against $\statorcurrent$ is not. Richter et al.
\cite{richter2014} distinguish magnetic-state matching from equality of
terminal currents in motoring/generating measurements.

For an algebraic illustration, take
$\flux=\matr L\magnetizingcurrent$ with
$\matr L=\operatorname{diag}(L_d,L_q)$ and
$\vect i_{\rm Fe}=c\rotationgenerator\flux$, $c=\omega_e/R_{\rm Fe}$.
Its dissipation in the same dq normalization is
$\|\omega_e\rotationgenerator\flux\|^2/R_{\rm Fe}>0$ for nonzero branch voltage;
physical three-phase power adds the $3/2$ factor under amplitude-invariant
coordinates. The storage law itself remains conservative. Eliminating its branch
current yields
\begin{equation}
 \frac{\partial\flux}{\partial\vect i_s}
 =\matr L(\matr I+c\rotationgenerator\matr L)^{-1}
 =\frac{1}{1+c^2L_dL_q}
 \begin{bmatrix}L_d&cL_dL_q\\-cL_dL_q&L_q\end{bmatrix}.
\end{equation}
Its terminal-current curl is $-2cL_dL_q/(1+c^2L_dL_q)$, despite conservative
storage magnetics. This illustrative constant branch is not identified for
the supplied machine. Indeed, its antisymmetric affine cross terms give
equal symmetry and curl equivalents; it does not automatically explain their
observed difference. Iron loss remains one open model component, not a
demonstrated cause.

Nor does simply relabeling currents transform a differential form correctly.
If $\magnetizingcurrent=f(\statorcurrent)$, then
\begin{equation}
 \nabla_{\vect i_s}\bigl[W'_{\rm mag}(f(\vect i_s))\bigr]
 =Df(\vect i_s)\trans\flux_{\rm mag}(f(\vect i_s)).
 \label{eq:current-pullback}
\end{equation}
The Jacobian transpose is required. This gradient is generally not the
unchanged flux vector plotted against new current labels. Magnetic coordinates,
their conjugate quantities, fixed state, and the subsystem being fitted must
therefore be specified together. Loss-aware synchronous-machine context is
distinct from induction-machine results; the Kullick--Hackl journal study
\cite{kullick2023} concerns induction machines and supplies no PMSM
co-energy validation for these terminal coordinates.

\subsection{Fitting inconsistent data is a model projection}
Symmetry and integrability are distinct tests. Parity can hold for a field
with nonzero circulation, while a conservative field can have a symmetry
axis different from the chosen coordinates. A small independent interpolation
error therefore says little about whether the field belongs to the expected
magnetic class. Conversely, a model constructed within that class can have
significant residuals if measurements are biased or the model is inappropriate.

For a fitted model write $\flux^{\rm rec}=\nabla\hat\coenergy+\vect r$.
This is a prediction--residual decomposition, not an assertion that the fitted
potential is the unknown true magnetic storage law in terminal coordinates.
For exactly mirrored current arguments at fixed state, the forbidden parity
channels of $\vect r$ equal those of the reconstructed field, because the
symmetric model has no contribution there. Requested-key pairing in the real
study does not ensure such exact measured-current geometry. The residual to the closest admissible map carries useful information,
but also includes basis limitations, regularization bias, and excluded
magnetics. The companion develops error hypotheses and identifiability rather
than automatically deleting this residual \cite{companion-diagnostics}.

The scalar potential presumes approximately quasi-static conservative
magnetics with fixed thermal and magnet state. Hysteresis, irreversible
demagnetization, dynamic iron-loss states, and strong frequency-dependent
effects require additional treatment. Real asymmetry can break reflection
without invalidating integrability. Inconsistency with this model is not proof
of a faulty measurement. The fitted field is the best representation in an
assumed class and chosen basis, not guaranteed ground truth.
For the real slices, a least-squares potential would project inconsistent
input onto the conservative class under its chosen support, weights, basis,
and regularization. A symmetry-constrained potential adds a second restriction.
Its zero curl and parity residuals are structural properties, not independent
evidence that discarded contributions were measurement errors or unphysical.
Raw reconstructed flux, model predictions, and their spatial residual must be
reported together, with high-current, global, and derivative metrics separated.
Winding resistance and inverter uncertainty are established identification
issues \cite{liu2018}, but that literature does not identify a cause in these
files. No real-data fitting success follows from the present admissibility test.

Scattered data require coverage in both current directions; boundary sparsity
can amplify derivative error. Symmetry reduces model freedom but does not
create independent observations on the unmeasured side. Extrapolation outside
the data support requires an explicit physical policy. Dimensionless scaling,
rank inspection, and region-based validation help expose fragile estimates.
Derivative quality should be assessed against independent references where
possible, not merely judged from a smooth plot.

The present synthetic study has a known representable polynomial, one random
seed, limited knot and weight choices, and a deliberately simple noise model.
RBFs are discussed conceptually but not numerically compared. Future evaluation
should vary seeds, nonrepresentable magnetic maps, noise correlations, outlier
patterns, gridded versus scattered data, boundary holes, and extrapolation.
Compare flux prediction, parity, reciprocity/curl, differential-inductance
smoothness and accuracy, robustness to bias, computational cost, and tuning
effort. No universal benchmark winner follows from a structural advantage.

\subsection{Preserve apparent tilt for joint model identification}
A visibly tilted reconstructed map can resemble the affine resistance
fingerprint without identifying a winding-resistance error. Automatic
untilting from symmetry alone is therefore unjustified. In a loss-aware
procedure, the same asymmetry may contain information about dissipative
currents, including iron loss, once magnetic states and conjugate coordinates
are matched. Erasing it before that analysis can remove an observable needed
to distinguish storage and loss. Neither the present residuals nor the
illustrative loss branch establish such a decomposition for these files.

For a sole constant resistance mismatch the reconstructed flux contribution
is proportional to $1/\omegae$. Multiple speeds test that hypothesis, but a
current-proportional voltage error has exactly the same scaling and remains
confounded. A promising next study would jointly model the storage potential,
loss states, and electrical measurement terms over speeds, independently
measured winding temperatures, and additional observables such as calibrated
terminal power, voltage, or magnetic-state references. The storage law must
use its declared conjugate coordinates and thermal state; loss and measurement
terms require their own observation relations. Common support, parameter-rank
checks, and held-out states are needed before interpreting a lower residual
as causal identification. This is a research direction, not an identified
multi-speed loss model or a real-data fitting result.
